\documentclass[11pt,a4paper]{article}
\usepackage[margin=1in]{geometry}
\usepackage{amsmath,amssymb,amsthm}
\usepackage{algorithm}
\usepackage{algpseudocode}
\usepackage{booktabs}
\usepackage{hyperref}

\theoremstyle{definition}
\newtheorem{definition}{Definition}[section]
\newtheorem{theorem}{Theorem}[section]
\newtheorem{lemma}[theorem]{Lemma}
\newtheorem{remark}{Remark}[section]

\title{\textbf{Rigorous Mathematical Specification, Residual Stream Variance Explosion, and Deep Residual Initialization (Fixup / GPT-2 Scaling)}}
\author{Team Scaibu \\ \small Email: \texttt{jaydeep.9664920749@gmail.com} \\ \small Architecture and Core Engine Team}
\date{October 2026}

\begin{document}
\maketitle

\begin{abstract}
This document presents the complete mathematical derivation and algorithmic specification of Deep Residual Initialization schemes (GPT-2 Residual Scaling, Zero-Init, and Fixup). In networks containing identity skip connections $x_{l+1} = x_l + F(x_l; W_l)$, if residual branches $F$ are initialized with standard unit variance, the variance of the residual stream grows linearly with depth ($\operatorname{Var}(x_L) \approx L \cdot \operatorname{Var}(x_0)$), causing gradient explosions and training instability in deep transformers (e.g., $L \ge 32$). Scaling the output projection layers by $\frac{1}{\sqrt{2L}}$ or initializing final branch weights to zero preserves strict isometric signal propagation at initialization. We formalize the residual variance accumulation theorem, present the complete algorithmic pseudo-code, derive complexity bounds, calculate a numerical example, and discuss architecture stability.
\end{abstract}

\section{Notation and Mathematical Preliminaries}

Let $x_l \in \mathbb{R}^d$ denote the residual stream embedding vector at layer $l \in \{1, \dots, L\}$, with transformation $x_{l+1} = x_l + F_l(x_l; W_l)$.

\begin{itemize}
    \item $L \in \mathbb{N}_{\ge 1}$: Total number of stacked residual blocks or transformer layers.
    \item $\sigma_0 > 0$: Unscaled baseline weight standard deviation (e.g., $\sigma_0 = 0.02$ or $1 / \sqrt{d_{\text{model}}}$).
    \item $s(L) \in \mathbb{R}_{\ge 0}$: Depth attenuation scale factor applied to residual branch output projections.
    \item $\sigma_{\text{scaled}} = s(L) \cdot \sigma_0$: Scaled initialization standard deviation.
    \item $\text{strategy} \in \{\text{scaled\_residual}, \text{zero\_init}, \text{fixup}\}$: Attenuation formulation.
\end{itemize}

\begin{table}[h]
\centering
\caption{Summary of Mathematical Symbols for Deep Residual Initialization}
\begin{tabular}{lll}
\toprule
\textbf{Symbol} & \textbf{Type / Domain} & \textbf{Description} \\
\midrule
$L$ & $\mathbb{N}_{\ge 1}$ & Total layer depth \\
$\sigma_0$ & $\mathbb{R}_{> 0}$ & Base standard deviation ($0.02$) \\
$s(L)$ & $\mathbb{R}_{\ge 0}$ & Multiplicative depth scale factor \\
$\sigma_{\text{scaled}}$ & $\mathbb{R}_{\ge 0}$ & Effective sampling standard deviation \\
\bottomrule
\end{tabular}
\end{table}

\section{Problem Definition and Core Concepts}

\subsection{Intuition and Motivation: The Residual Stream Variance Blowup}
In a pre-layer-norm Transformer or ResNet without attenuation, the residual recurrence relation is:
\begin{equation}
x_L = x_0 + \sum_{l=1}^L F_l(x_l)
\end{equation}
Assuming independent zero-mean residual branches with variance $\operatorname{Var}(F_l(x)) \approx \sigma_F^2$, the variance of the residual stream expands as:
\begin{equation}
\operatorname{Var}(x_L) = \operatorname{Var}(x_0) + \sum_{l=1}^L \operatorname{Var}(F_l(x_l)) \approx \operatorname{Var}(x_0) + L \sigma_F^2
\end{equation}
For deep models ($L = 96$ in GPT-3 175B), activations at layer 96 are $100\times$ larger than at layer 1, causing attention softmax distributions to become nearly one-hot delta functions, collapsing gradient propagation. Attenuating the output projection weights of each residual block by $s(L) = \frac{1}{\sqrt{2L}}$ (accounting for 2 sub-layers per block: MHA and MLP) bounds total variance growth to $\mathcal{O}(1)$.

\subsection{Formal Mathematical Definition}
\begin{definition}[Residual Scaling Schemes]
Given total layer count $L \ge 1$ and baseline standard deviation $\sigma_0 > 0$:
\begin{enumerate}
    \item \textbf{Scaled Residual (GPT-2 / Radford et al., 2019)}:
    \begin{equation}
    s(L) = \frac{1}{\sqrt{2L}}, \quad \sigma_{\text{scaled}} = \frac{\sigma_0}{\sqrt{2L}} \label{eq:gpt2_scale}
    \end{equation}
    \item \textbf{Zero-Init (Goyal et al., 2017; Bachlechner et al., 2021)}:
    \begin{equation}
    s(L) = 0.0, \quad \sigma_{\text{scaled}} = 0.0 \label{eq:zero_init}
    \end{equation}
    (Initializes the last linear projection of each residual block to identically zero, so $F(x) = \mathbf{0}$ and $x_{l+1} \equiv x_l$ at step $0$).
    \item \textbf{Fixup (Zhang et al., 2019)}:
    \begin{equation}
    s(L) = L^{-1/4}, \quad \sigma_{\text{scaled}} = \sigma_0 \cdot L^{-1/4} \label{eq:fixup_scale}
    \end{equation}
\end{enumerate}
\end{definition}

\section{Algorithmic Specification}

The computational pipeline implemented in \texttt{NnAlgoDeepResidualInit} is detailed below.

\begin{algorithm}[H]
\caption{Deep Residual Initialization Scale Factor Computation}
\begin{algorithmic}[1]
\Require Total layer count $L \ge 1$, base standard deviation $\sigma_0 > 0$
\Require Strategy $\in \{\text{scaled\_residual}, \text{zero\_init}, \text{fixup}\}$
\Ensure Scaled standard deviation $\sigma_{\text{scaled}}$, attenuation factor $s(L)$
\If{$L < 1$ \textbf{or} $\sigma_0 \le 0$}
    \State \textbf{throw} PreconditionFailureException(``Layer count must be positive and sigma must be > 0'')
\EndIf
\If{$\text{strategy} = \text{scaled\_residual}$}
    \State $s \gets 1.0 / \sqrt{2.0 \cdot \text{float}(L)}$
    \State $\sigma_{\text{scaled}} \gets \sigma_0 \cdot s$
\ElsIf{$\text{strategy} = \text{zero\_init}$}
    \State $s \gets 0.0$
    \State $\sigma_{\text{scaled}} \gets 0.0$
\ElsIf{$\text{strategy} = \text{fixup}$}
    \State $s \gets (\text{float}(L))^{-0.25}$
    \State $\sigma_{\text{scaled}} \gets \sigma_0 \cdot s$
\Else
    \State \textbf{throw} PreconditionFailureException(``Unrecognized initialization strategy'')
\EndIf
\State \Return $\{\text{scaled\_std}: \sigma_{\text{scaled}}, \text{scale\_factor}: s\}$
\end{algorithmic}
\end{algorithm}

\section{Theoretical Guarantees and Isometric Signal Propagation}

\begin{theorem}[Variance Invariance of Attenuated Residual Stream]
Let each transformer block contain 2 residual branches (Self-Attention and Feed-Forward) with linear output projections $W_{\text{proj}}^{(l)} \sim \mathcal{N}\left(0, \frac{\sigma_0^2}{2L}\right)$. Then the variance of the residual stream at output layer $L$ satisfies:
\begin{equation}
\operatorname{Var}(x_L) = \operatorname{Var}(x_0) + \mathcal{O}(1)
\end{equation}
independent of total depth $L$.
\end{theorem}

\begin{proof}
Let $x_{l, 1} = x_l + F_{\text{attn}}(x_l)$ and $x_{l+1} = x_{l, 1} + F_{\text{mlp}}(x_{l, 1})$.
The incremental variance added across both sub-layers in block $l$ is:
\begin{equation}
\Delta \operatorname{Var}_l = \operatorname{Var}(F_{\text{attn}}) + \operatorname{Var}(F_{\text{mlp}}) \approx 2 \cdot \left(\frac{\sigma_0^2}{2L}\right) \operatorname{Var}(x_0) = \frac{\sigma_0^2}{L} \operatorname{Var}(x_0)
\end{equation}
Summing over all $L$ layers:
\begin{equation}
\operatorname{Var}(x_L) = \operatorname{Var}(x_0) + \sum_{l=1}^L \Delta \operatorname{Var}_l = \operatorname{Var}(x_0) + L \left(\frac{\sigma_0^2}{L} \operatorname{Var}(x_0)\right) = \operatorname{Var}(x_0)(1 + \sigma_0^2)
\end{equation}
which is strictly $\mathcal{O}(1)$ and uniformly bounded for all $L \to \infty$.
\end{proof}

\subsection{Limitations}
\begin{itemize}
    \item \textbf{Residual Branch Selective Application}: Attenuation must ONLY be applied to the residual branch projection matrices ($W_{\text{out}}$ in MHA, $W_{\text{down}}$ in MLP/FFN), NOT to the initial embeddings or feed-forward expansion layers.
\end{itemize}

\section{Complexity Analysis}

\subsection{Time and Space Complexity}
\begin{itemize}
    \item \textbf{Time}: Exactly $\mathcal{O}(1)$ scalar arithmetic operations.
    \item \textbf{Space}: $\mathcal{O}(1)$ memory.
\end{itemize}

\section{Worked Numerical Example}

Consider a $L = 32$-layer Transformer with $\sigma_0 = 0.02$.

\begin{enumerate}
    \item \textbf{Scaled Residual (GPT-2)}:
    \begin{align}
    s &= \frac{1}{\sqrt{2 \times 32}} = \frac{1}{\sqrt{64}} = \frac{1}{8.0} = 0.1250 \\
    \sigma_{\text{scaled}} &= 0.02 \times 0.1250 = 0.00250
    \end{align}
    \item \textbf{Fixup}:
    \begin{align}
    s &= 32^{-0.25} \approx 0.4204482 \\
    \sigma_{\text{scaled}} &= 0.02 \times 0.4204482 \approx 0.00840896
    \end{align}
    \item \textbf{Zero-Init}:
    \begin{equation}
    s = 0.0, \quad \sigma_{\text{scaled}} = 0.0
    \end{equation}
\end{enumerate}

\section{Numerical Considerations and Edge Cases}

\begin{itemize}
    \item \textbf{Zero-Init Implementation}: Initializing final linear weights to zero ensures exact identity mapping $x_{l+1} \equiv x_l$ at step $0$, completely bypassing any signal degradation during the very first forward/backward pass.
    \item \textbf{LayerNorm Interaction}: In Pre-LN Transformers, combining Pre-LN with $1/\sqrt{2L}$ residual scaling prevents gradient vanishing at initial layers.
\end{itemize}

\begin{thebibliography}{9}
\bibitem{radford2019} A.~Radford et al., ``Language models are unsupervised multitask learners,'' \emph{OpenAI Blog}, vol.~1, no.~8, p.~9, 2019.
\bibitem{zhang2019fixup} H.~Zhang, Y.~N.~Dauphin, and T.~Ma, ``Fixup initialization: Residual learning without normalization,'' in \emph{International Conference on Learning Representations (ICLR)}, 2019.
\bibitem{bachlechner2021} T.~Bachlechner et al., ``ReZero is all you need: Fast convergence at large depth,'' in \emph{Conference on Uncertainty in Artificial Intelligence (UAI)}, pp.~1352--1361, 2021.
\end{thebibliography}

\end{document}
